\rtmtight
\begin{tblr}{width=\textwidth,colspec={|Q[c,m,2.1cm]|X[l,m]|},hlines,vlines,hspan=minimal,rowsep=1.5pt,colsep=3pt,row{1}={bg=headgray,font=\bfseries}}
\SetCell{c,m}\hfil\logo\hfil & \textbf{POLITEKNIK NEGERI MALANG}\par\textbf{JURUSAN TEKNOLOGI INFORMASI}\par\textbf{PROGRAM STUDI: D4 TEKNIK INFORMATIKA} \\
\SetCell[c=2]{l} \textbf{RENCANA TUGAS MAHASISWA (RTM) DAN RUBRIK PENILAIAN} & \\
Nama Mata Kuliah & Praktikum Pemrograman Berbasis Objek \\
Kode Mata Kuliah & RTI253008 \quad \textbf{sks / Jam perminggu:} 2 SKS/4 jam \quad \textbf{Semester:} 3 \\
Dosen Pengampu & Tim Pengajar Program Studi D4 TI \\
\end{tblr}
\assessmentblock{Tugas Jobsheet Praktikum PBO}{Case Method dan praktik pemrograman Java}{SCPMK0210-02601, SCPMK0704-02602, SCPMK0903-02603}{Mahasiswa menyelesaikan jobsheet praktikum sesuai topik mingguan mencakup OOP dasar, OOP lanjut, GUI Swing, dan koneksi database JDBC.}{Menganalisis kebutuhan kasus pada jobsheet, mengimplementasikan kode Java, menguji hasil eksekusi, mendokumentasikan laporan, dan mengumpulkan evidence.}{Kode sumber Java, laporan jobsheet, screenshot eksekusi, dan/atau bahan presentasi.}{Ketepatan implementasi konsep OOP & 12 & Kelas, enkapsulasi, inheritance, dan relasi diimplementasikan dengan benar sesuai spesifikasi jobsheet tanpa error runtime. & Implementasi sebagian besar benar, tetapi ada kekurangan pada penerapan enkapsulasi atau relasi kelas. & Implementasi tidak menerapkan konsep OOP yang diminta atau menghasilkan error runtime. \\ \\
Pengerjaan fitur OOP lanjut: polymorfisme, interface, dan koleksi & 10 & Fitur OOP lanjut diimplementasikan dengan benar, semua kasus uji jobsheet lulus, dan hasilnya terdokumentasi. & Sebagian besar fitur lanjut diimplementasikan tetapi ada kasus uji yang tidak lulus. & Fitur lanjut tidak diimplementasikan atau tidak berfungsi. \\ \\
Kualitas aplikasi GUI dan JDBC & 8 & GUI Swing berjalan responsif, koneksi JDBC berhasil, dan operasi CRUD berfungsi sesuai spesifikasi. & GUI dan JDBC berjalan untuk fitur utama tetapi beberapa operasi belum lengkap. & GUI tidak berjalan atau koneksi JDBC gagal. \\}

\assessmentblock{Kuis Konsep OOP}{Kuis}{SCPMK0210-02601, SCPMK0704-02602}{Kuis tertulis untuk mengukur pemahaman konsep OOP dasar (Kuis 1) dan OOP lanjut (Kuis 2) menggunakan bahasa Java.}{Menjawab soal pilihan ganda dan/atau uraian terkait konsep OOP, kelas, enkapsulasi, inheritance, overriding, interface, abstrak, dan polymorfisme.}{Jawaban tertulis kuis.}{Penguasaan konsep OOP dasar & 3 & Konsep kelas, enkapsulasi, inheritance, dan relasi kelas dijelaskan dengan benar dan contoh penerapannya tepat dalam kode Java. & Sebagian besar konsep dipahami tetapi penjelasan atau contoh masih kurang lengkap. & Konsep OOP dasar tidak dipahami atau penjelasan keliru. \\ \\
Kemampuan analisis kode Java & 4 & Kode Java dianalisis dengan tepat, output diprediksi benar, dan error diidentifikasi secara akurat. & Sebagian besar analisis kode benar tetapi beberapa prediksi output atau identifikasi error masih keliru. & Analisis kode tidak tepat atau tidak menunjukkan pemahaman sintaks Java. \\ \\
Penguasaan konsep OOP lanjut & 3 & Overriding, overloading, interface, kelas abstrak, dan polymorfisme dijelaskan dan diterapkan dengan benar pada kasus. & Sebagian konsep OOP lanjut dipahami tetapi penerapannya pada kasus belum lengkap. & Konsep OOP lanjut tidak dipahami atau penerapannya keliru. \\}

\assessmentblock{UTS Praktikum OOP Dasar}{UTS berbasis praktikum}{SCPMK0210-02601, SCPMK0704-02602}{Evaluasi tengah semester untuk mengukur kemampuan implementasi konsep OOP dasar dan lanjut menggunakan Java secara mandiri.}{Mengimplementasikan program Java yang menerapkan kelas, enkapsulasi, inheritance, relasi kelas, dan elemen OOP lanjut sesuai soal ujian.}{Kode sumber Java yang berjalan, laporan singkat, dan screenshot eksekusi.}{Implementasi kelas dan enkapsulasi & 5 & Kelas diimplementasikan dengan atribut private, getter/setter yang benar, konstruktor yang sesuai, dan relasi antarkelas terpenuhi. & Implementasi kelas ada tetapi penggunaan modifier, konstruktor, atau relasi belum sepenuhnya tepat. & Implementasi tidak menerapkan enkapsulasi atau struktur kelas tidak benar. \\ \\
Implementasi inheritance dan overriding & 5 & Inheritance diimplementasikan dengan benar menggunakan extends, super, dan override sesuai hierarki kelas yang dirancang. & Inheritance ada tetapi penggunaan super, override, atau hierarki kelas masih belum optimal. & Inheritance tidak diterapkan atau implementasinya menghasilkan error runtime. \\ \\
Kualitas kode dan dokumentasi & 5 & Kode bersih, terdokumentasi dengan komentar yang memadai, nama variabel/metode deskriptif, dan mengikuti konvensi Java. & Kode berfungsi tetapi dokumentasi atau konvensi penamaan masih belum konsisten. & Kode sulit dibaca, tidak terdokumentasi, atau tidak mengikuti konvensi Java. \\}

\assessmentblock{PBL Aplikasi Java GUI dan Database}{Project Based Learning}{SCPMK0704-02602, SCPMK0903-02603}{Mahasiswa mengembangkan aplikasi Java lengkap berbasis GUI Swing dan database menggunakan seluruh prinsip OOP yang telah dipelajari dalam satu proyek terintegrasi.}{Merancang arsitektur OOP, mengimplementasikan GUI Swing dan JDBC, mengintegrasikan seluruh fitur, menguji fungsionalitas, mendokumentasikan, dan mempresentasikan.}{Aplikasi Java yang berjalan, kode sumber, diagram kelas, laporan pengujian, dan/atau bahan presentasi.}{Arsitektur OOP dan penggunaan polymorfisme & 8 & Arsitektur OOP diterapkan konsisten, polymorfisme diimplementasikan melalui interface atau kelas abstrak yang tepat, dan desain kelas terdokumentasi. & Arsitektur OOP ada tetapi penggunaan polymorfisme atau dokumentasi desain belum konsisten. & Arsitektur tidak berbasis OOP atau polymorfisme tidak diterapkan. \\ \\
Fungsionalitas GUI dan koneksi database & 8 & GUI Swing lengkap dengan event handling yang benar, JDBC terhubung, dan operasi CRUD berjalan sesuai spesifikasi. & Sebagian besar fungsionalitas berjalan tetapi beberapa fitur GUI atau operasi database belum lengkap. & GUI tidak lengkap atau koneksi database gagal. \\ \\
Kualitas kode, pengujian, dan presentasi & 6 & Kode terstruktur, skenario uji terdokumentasi, dan presentasi menjelaskan arsitektur serta keputusan implementasi secara meyakinkan. & Kode berfungsi dan presentasi cukup jelas, tetapi pengujian atau dokumentasi belum lengkap. & Kode tidak terstruktur, tidak diuji, atau presentasi tidak menjelaskan arsitektur desain. \\}

\assessmentblock{UAS Praktikum: Demo Aplikasi Final}{UAS berbasis demo dan defense}{SCPMK0903-02603}{Mahasiswa mendemonstrasikan aplikasi Java final yang dikembangkan dan mempertahankan keputusan implementasi OOP kepada dosen penguji.}{Mempersiapkan aplikasi berjalan, mempresentasikan fitur, mendemonstrasikan operasi, dan menjawab pertanyaan terkait keputusan desain OOP.}{Aplikasi Java yang berjalan, kode sumber, dokumentasi teknis, dan laporan akhir.}{Kelengkapan dan fungsionalitas aplikasi & 8 & Semua fitur yang diminta berjalan, GUI responsif, koneksi database berfungsi, dan tidak ada error kritis saat demo. & Fitur utama berjalan tetapi beberapa fungsi atau koneksi database masih ada kekurangan minor. & Banyak fitur tidak berfungsi atau aplikasi tidak dapat didemonstrasikan. \\ \\
Penerapan prinsip OOP secara menyeluruh & 8 & Seluruh prinsip OOP (enkapsulasi, inheritance, polymorfisme, abstraksi) diterapkan secara konsisten dan dapat ditelusuri dalam kode. & Sebagian besar prinsip OOP diterapkan tetapi beberapa bagian tidak konsisten atau kurang optimal. & Prinsip OOP tidak diterapkan atau aplikasi menggunakan pendekatan prosedural. \\ \\
Defense dan kualitas dokumentasi & 7 & Mahasiswa mampu menjelaskan dan mempertahankan setiap keputusan desain OOP dengan evidence dari kode, dan dokumentasi lengkap tersedia. & Defense cukup jelas tetapi beberapa keputusan tidak dapat dipertahankan dengan baik atau dokumentasi belum lengkap. & Tidak mampu menjelaskan desain OOP atau dokumentasi tidak tersedia. \\}

\begin{tblr}{width=\textwidth,colspec={|Q[l,m,3.2cm]|X[l,m]|},hlines,vlines,hspan=minimal,row{1-3}={bg=headgray},rowsep=1.5pt,colsep=3pt}
Jadwal Pelaksanaan: & Minggu 1--17 sesuai RPS. Setiap evaluasi memiliki RTM dan rubrik multi-kriteria. \\
Lain-Lain Yang Diperlukan: & LMS, repositori/kode sumber Java, IDE, perangkat lunak sesuai mata kuliah, dan perangkat presentasi. \\
Pustaka: & Mengacu pustaka utama dan pendukung pada bagian RPS. \\
\end{tblr}
